\chapter{Introduction}
\label{ch:introduction}

Modern pharmacotherapy requires the integration of well-established pharmacological 
knowledge across multiple biological and clinical levels. Medical knowledge graphs provide a 
structured representation of this knowledge, capturing the complex mechanisms through 
which drugs interact with the human organism, with one another, and with intermediate 
biological processes. These interactions may propagate through multiple mechanistic pathways 
and ultimately converge on a observable clinical endpoint, particularly in the context of 
suboptimal or complex pharmacotherapy.

The topology of a knowledge graph describes which pathways and relationships are considered plausible or
mechanistically meaningful based on literature, but it does not by itself determine
how strongly these relationships are supported by observations in a particular population,
nor how the represented mechanisms are expressed in an individual patient.
On the other hand, patient-level data provide empirical information on drug exposure, laboratory 
outcomes, and diagnoses, but these observations are often incomplete, noisy, and 
affected by selection artifacts. As a result, the observed data do not directly reveal the 
underlying mechanistic structure linking drug exposures to intermediate processes and, ultimately, 
to clinical endpoints.
This creates a natural methodological distinction between two complementary sources of information
-- graph-structured knowledge and patient-level empirical evidence.

The central aim of this thesis is to investigate how these two sources can be integrated within
a unified framework so that they interact synergistically.
Based on this aim, we formulate the following theses:

\paragraph{Thesis 1.}
Empirical patient-level data, when integrated with the structural information encoded in a 
knowledge graph, can improve the prediction of missing directed links.

\paragraph{Thesis 2.}
The directed causal graph can serve as a shared scaffold for predicting multiple
clinical events in individual patients.

\medskip

As we combine graph-based structure and empirical data, this leads us to use
graph neural networks (GNNs) for prediction.
As the signal in GNNs can be averaged out, we explore
methods for incorporating additional signal into GNNs through associative statistics.
Inspired by the Wide \& Deep framework proposed by Cheng et al.~\cite{cheng2016wide}, 
we formulate the third thesis:

\paragraph{Thesis 3.}
Explicit empirical associations can provide additional signal to GNNs and improve prediction performance of GNNs.

\medskip

%A medical digital twin is here understood as an updatable, patient-specific representation
%of the mechanisms, exposures and observations relevant to a clinical decision,
%rather than as a complete virtual replica of a person.
The importance of these theses is reflected in the current research on
patient-specific modelling and digital twins (digital representations of patients).
The direction proposed is discussed in work on causal graph neural networks in healthcare,
which connects mechanistic graph structure with patient-specific data.
That literature also notes that models relying predominantly on statistical associations
may be sensitive to changes between populations, institutions, and data-generating
environments, and that causal or mechanistic structure is hypothesized to yield
representations that are more stable under such changes---while emphasizing remaining
challenges related to distribution shifts, validation, computational complexity, and
the distinction between causal claims and statistical
associations~\cite{mesinovic2025causal}.

% A related idea appears more broadly in machine learning, where explicit statistical
% representations and learned neural representations do not
% necessarily have to compete with each other.
% The Wide \& Deep framework introduced by Cheng et al.\ combines a wide component, designed to 
% exploit observed feature interactions, with a deep neural component that generalizes to previously 
% unseen feature combinations.
% Although developed for recommender systems, it provides an important methodological example
% of jointly learning from explicit empirical relationships and learned representations
% instead of relying exclusively on increasingly complex neural architectures~\cite{cheng2016wide}.
% This combination of explicit statistics with a neural encoder is the setting of Thesis~3.

% In this thesis the associational component is a Hierarchical Correlation Reconstruction
% (HCR) representation of empirical dependence.

Based on that and taking into consideration formulated theses,
there are defined two complementary prediction tasks.

\paragraph{Task~A.}
Patient-level observations are used together with the available graph structure to
evaluate candidate directed relations and perform link prediction.
%The objective is to determine whether an unobserved or withheld relation is compatible both
%with its structural position in the heterogeneous graph and with statistical evidence derived
%from the corresponding patient variables.

\paragraph{Task~B.}
Patient-level clinical outcome prediction based on the graph structure.
The graph structure is treated as a shared mechanistic scaffold, while patient-specific observations are
used to predict clinical endpoints for individual patients.
%Opposite to Task~A, Task~B asks whether 
%the graph and the available patient information together support the prediction of a particular clinical outcome.
\medskip

Together, the two tasks can be viewed as complementary parts of a broader iterative framework.
In real-world applications, patient-level evidence can be used systematically to evaluate and 
extend the structured knowledge representation, while the resulting graph can in turn provide context for 
patient-level prediction.
In this way, the knowledge graph is not treated only as a static repository of prior knowledge, 
but as a structured representation that can, in principle, be confronted with accumulating empirical evidence 
at scale.
% \begin{quote}
% \textit{mechanistic knowledge} $\rightarrow$ \textit{patient prediction} $\rightarrow$ \textit{empirical evidence} $\rightarrow$ \textit{evaluation/enrichment of knowledge}.
% \end{quote}
% This loop is the intended research direction. The experiments test both of its legs;
% they do not implement a closed, self-updating cycle.

\begin{figure}[H]
    \centering
    \includegraphics[width=0.92\textwidth,height=0.32\textheight,keepaspectratio]{fig_loop_task_a_b.png}
    \caption{Intended research direction: an iterative loop between Task~A
    (evaluation of graph relations) and Task~B (patient-level endpoint prediction).}
    \label{fig:loop}
    \label{fig:intro-task-loop}
\end{figure}